


I.2 Méthodes
I.2.1) Mortalité routière
i. Détection de carcasses à pied et à partir d'un véhicule
Nous avons comparé les détections de carcasses et d'individus vivants de tous les taxons (amphibiens, reptiles, mammifères, oiseaux) sur la chaussée de l'A10 et la R104 en utilisant le suivi à pied par jumelles, et le suivi avec caméra réalisé par automobile. Lors d'une même journée d'échantillonnage, le même segment de route était parcouru à l'aide des deux méthodes : suivi à pied par jumelles entre 15 et 30 m de la chaussée et suivi avec caméra fixée sur un véhicule circulant sur la chaussée. Afin de réduire la probabilité de retrait des carcasses sur la route par des prédateurs ou des automobiles, ces suivis étaient effectués dès le lever du soleil. Lors d'une journée donnée, nous avons déterminé aléatoirement l’ordre des méthodes utilisées. Le but était de comparer l'efficacité des deux méthodes à détecter sur la route les mêmes carcasses et les mêmes individus vivants. Il était donc important de réaliser les deux suivis dans un court laps de temps (même matinée). L’utilisation des deux méthodes permettrait de mieux estimer le nombre réel de carcasses sur la route.
Le suivi par jumelles consistait à parcourir à pied les sentiers situés à proximité des tronçons de l'A10 et de la R104, en restant entre 15 et 30 m de la chaussée. Au moins deux observateurs réalisaient l'inventaire. Le port d'un gilet de sécurité réfléchissant et d'un casque était obligatoire en tout temps pour augmenter la visibilité du personnel sur le terrain. Les jumelles utilisées étaient des Nikon Prostaff P3 8X42 (Nikon, Tokyo, Japon). Toutes les carcasses observées sur les routes étaient géolocalisées à l'aide d'un GPS GARMIN GPSMAP 79s (Garmin, Schaffhausen, Suisse). Le suivi par caméra réalisé par voiture utilisait une caméra GoPro Hero Black 11 (GoPro, San Mateo, Californie, États-Unis) afin d'enregistrer les images sur la chaussée avec une qualité 4K, un ratio 4:3 et 30 images par seconde. Lors d'une session d'échantillonnage, la caméra GoPro était insérée dans un étui imperméable transparent et fixée à un trépied magnétique (3-Footed Monster, Sydney, Australie). Le trépied magnétique était fixé au milieu du capot avant de la voiture (Figure 3). La caméra pouvait être contrôlée à distance par protocole Bluetooth à l'aide de l'application gratuite Quik pour Android. Nous lancions l'enregistrement lorsque le véhicule entrait dans le segment suivi jusqu'à la fin du parcours. Les fichiers vidéo ont été visionnés et les données d'observation sur la route ont été saisies en notant le moment (minutes, secondes) et la position (coordonnées géographiques) de la séquence vidéo.